\section*{Hoofdstuk \ref{ch:b3:lagrangian-mechanics} --- Lagrangiaanse mechanica}

\begin{solution}{exo:b3:lagrangian-mechanics:1}
(a) 2 (twee hoeken op het oppervlak van de schaal). (b) 1: de abscis $x$
langs de helling, want de draaihoek ligt door het rollen aan haar vast,
$x =
R\phi$. (c) 3: $X$, $\theta_1$, $\theta_2$. (d) 5: drie voor het
middelpunt, twee voor de richting van de staaf. (e) 2: elk één hoek. De
rolvoorwaarde is een betrekking tussen \emph{snelheden}, $\dot x =
R\dot\phi$; in dit vlakke vraagstuk is zij meteen te integreren tot $x = R\phi +
\text{const}$, en dus holonoom. (Voor een bal die op een vlak
rolt is zij niet integreerbaar, en heeft de lagrangiaanse mechanica een
uitbreiding nodig.)
\end{solution}

\begin{solution}{exo:b3:lagrangian-mechanics:2}
(a) $[L] = \unit{J}$; $[p_\theta] = \unit{kg.m^2/s}$ --- een
impulsmoment, want $\theta$ is dimensieloos. (b) $p_\theta =
m\ell^2\dot\theta$ is het impulsmoment van het gewicht om het draaipunt.
(c) $\ddot\theta = -(g/\ell)\theta$ geeft $T = 2\pi\sqrt{\ell/g}$. (d)
Met $E_p$ gemeten vanaf het evenwicht heeft een harmonische trilling
gelijke tijdgemiddelde kinetische en potentiële energie, dus $\langle
L\rangle = 0$ en $S = 0$ over een hele periode --- de berekening
bevestigt het: $S = \int_0^T\tfrac12 m\ell^2\theta_0^2[\omega^2\sin^2
\omega t - (g/\ell)\cos^2\omega t]\,\dd t = 0$ omdat $\omega^2 =
g/\ell$.
\end{solution}

\begin{solution}{exo:b3:lagrangian-mechanics:3}
(a) Zij $x$ de daling van $m_1$ (zodat $m_2$ met $x$ stijgt): $L =
\tfrac12(m_1 + m_2)\dot x^2 + (m_1 - m_2)gx$. (b) $(m_1 + m_2)\ddot x =
(m_1 - m_2)g$: $a = (m_1 - m_2)g/(m_1 + m_2)$. (c) De spankracht grijpt
aan beide uiteinden van een onrekbaar koord aan: in elke toegelaten
verplaatsing heffen haar arbeiden elkaar op, dus komt zij nooit in $L$
voor. (d) Newton op $m_1$ alleen: $T =
m_1(g - a) = 2m_1m_2g/(m_1 + m_2)$.
\end{solution}

\begin{solution}{exo:b3:lagrangian-mechanics:4}
(a) $\varphi$ en $z$: $p_\varphi = mr^2\dot\varphi$ (het impulsmoment om
de as) en $p_z = m\dot z$ blijven behouden. (b) $r$ komt in $L$ voor via
$r^2\dot\varphi^2$ en is dus niet cyclisch: $\dot p_r =
mr\dot\varphi^2 \neq 0$. Daar is niets mis mee --- $p_r = m\dot r$ is de
radiale component van een constante vector $\vect p$, en een component
langs een \emph{draaiende} richting hoeft niet constant te zijn. (c)
$m\ddot r =
mr\dot\varphi^2$: de centrifugale term, de prijs van
coördinaten die aan draaiende richtingen vastzitten. (d) Voor een lijn
op afstand $b$ met snelheid $v$ doorlopen: $r = \sqrt{b^2 + v^2t^2}$,
$r^2\dot\varphi = bv$; dan is $\ddot r = b^2v^2/r^3 = r(bv/r^2)^2 = r\dot\varphi^2$.
\end{solution}

\begin{solution}{exo:b3:lagrangian-mechanics:5}
(a) Positie van het blok $(X + s\cos\alpha,\ -s\sin\alpha)$, dus
\[
L = \tfrac12 M\dot X^2 + \tfrac12 m\big(\dot X^2 +
2\dot X\dot s\cos\alpha + \dot s^2\big) + mgs\sin\alpha .
\]
(b) $X$ is cyclisch: $P = (M + m)\dot X + m\dot s\cos\alpha$ blijft
behouden --- de totale horizontale impuls, want er werkt geen uitwendige
horizontale kracht. (c) De vergelijking voor $s$ luidt $\ddot s + \ddot
X\cos\alpha = g\sin\alpha$; met $\ddot X = -m\ddot s\cos\alpha/(M +
m)$
uit (b),
\[
\ddot s = \frac{(M + m)g\sin\alpha}{M + m\sin^2\alpha} , \qquad
\ddot X = -\frac{mg\sin\alpha\cos\alpha}{M + m\sin^2\alpha} .
\]
(d) $M \to \infty$: $\ddot s \to g\sin\alpha$, het vaste hellend vlak;
$\alpha \to 90^\circ$: $\ddot s \to g$, $\ddot X \to 0$ --- vrije val
langs een verticale wand, zonder duw tegen de wig.
\end{solution}

\begin{solution}{exo:b3:lagrangian-mechanics:6}
(a) $L = \tfrac12 m\ell^2(\dot\theta^2 + \sin^2\theta\,\dot\varphi^2) +
mg\ell\cos\theta$. (b) $\varphi$ is cyclisch: $p_\varphi =
m\ell^2\sin^2\theta\,\dot\varphi$, het verticale impulsmoment. (c)
Na eliminatie van $\dot\varphi$ blijft $\tfrac12 m\ell^2\dot\theta^2 +
U_{\text{eff}}(\theta)$ behouden, met
\[
U_{\text{eff}}(\theta) = \frac{p_\varphi^2}{2m\ell^2\sin^2\theta}
 - mg\ell\cos\theta :
\]
een wand bij $\theta = 0$ en $\theta = \pi$ (voor $p_\varphi \neq 0$)
met één minimum ertussen --- het gewicht nutateert tussen twee cirkels.
(d) In het minimum is $\dot\theta = 0$ met $\theta = \theta_0$ constant:
$p_\varphi^2\cos\theta_0/m\ell^2\sin^3\theta_0 = mg\ell\sin\theta_0$;
invullen van $p_\varphi = m\ell^2\sin^2\theta_0\,\omega$ geeft
$\ell\omega^2\cos\theta_0 = g$.
\end{solution}

\begin{solution}{exo:b3:lagrangian-mechanics:7}
(a) Gewicht in $(X + \ell\sin\theta,\ -\ell\cos\theta)$:
\[
L = \tfrac12(M + m)\dot X^2 + m\ell\cos\theta\,\dot X\dot\theta
 + \tfrac12 m\ell^2\dot\theta^2 + mg\ell\cos\theta .
\]
(b) $X$ is cyclisch: $(M + m)\dot X + m\ell\cos\theta\,\dot\theta$
blijft behouden --- de horizontale impuls van het hele stelsel (de rail
duwt alleen verticaal). (c) Kleine hoeken: $(M + m)\ddot X +
m\ell\ddot\theta = 0$ en $\ell\ddot\theta + \ddot X + g\theta = 0$; na
eliminatie van $\ddot X$ is $\ell\ddot\theta\,[1 - m/(M + m)] = -g\theta$, dus $\Omega^2 = (M + m)g/M\ell = (1 + m/M)\,g/\ell$. (d) $M \to \infty$: het vaste draaipunt, $\Omega^2 = g/\ell$. Kleine $M$: het
wagentje wijkt tegengesteld aan het gewicht terug, de slingerbeweging
verloopt om een punt tussen beide (het vaste massamiddelpunt), wat de
effectieve slinger verkort --- vandaar de hogere frequentie, die
divergeert als $M \to 0$.
\end{solution}

\begin{solution}{exo:b3:lagrangian-mechanics:8}
(a) $\vect A = \tfrac12 B(-y, x, 0)$:
$L = \tfrac12 m(\dot x^2 + \dot y^2 + \dot z^2) + \tfrac12 qB(x\dot y -
y\dot x)$. (b) De vergelijking voor $x$: $\dd(m\dot x - \tfrac12 qBy)/\dd t =
\tfrac12 qB\dot y$, dat wil zeggen $m\ddot x = qB\dot y$; evenzo
$m\ddot y =
-qB\dot x$: cirkelbeweging bij $\omega_{\text{c}} = qB/m$,
met $\dot z$ constant. (c) $p_x = m\dot x - \tfrac12 qBy \neq m\dot x$;
noch $x$ noch $y$ is cyclisch ($\partial L/\partial x = \tfrac12 qB\dot y$), dus blijft geen van beide impulsen behouden --- alleen combinaties
als $m\dot x -
qBy$ (ga na dat de afgeleide ervan verdwijnt). (d) In
cilindercoördinaten is $A_\varphi = \tfrac12 Br$:
$L = \tfrac12 m(\dot r^2 + r^2\dot\varphi^2 + \dot z^2) + \tfrac12
qBr^2\dot\varphi$; $\varphi$ is cyclisch, $p_\varphi = mr^2\dot\varphi +
\tfrac12 qBr^2$. Op een cirkel met straal $R$ rond de as is
$\dot\varphi = -\omega_{\text{c}}$, dus $p_\varphi = -qBR^2 + \tfrac12
qBR^2 = -\tfrac12 qBR^2$.
\end{solution}

\begin{solution}{exo:b3:lagrangian-mechanics:9}
(a) $h = \dot\theta\,\partial L/\partial\dot\theta - L = \tfrac12
mR^2\dot\theta^2 - \tfrac12 m\omega^2R^2\sin^2\theta - mgR\cos\theta =
\tfrac12 mR^2\dot\theta^2 + U_{\text{eff}}$; $\dd h/\dd t =
\dot\theta\,[mR^2\ddot\theta + U_{\text{eff}}'(\theta)] = 0$ volgens de
bewegingsvergelijking. (b) $E = \tfrac12 mR^2\dot\theta^2 + \tfrac12
m\omega^2R^2\sin^2\theta - mgR\cos\theta = h +
m\omega^2R^2\sin^2\theta$. (c) $P = \dd E/\dd t =
m\omega^2R^2\sin2\theta\,\dot\theta$: positief zolang de kraal van de as
weg klimt (de motor werkt tegen de traagheid van de kraal in), negatief
op de terugweg. (d) $U_{\text{eff}}'' = mgR\cos\theta -
m\omega^2R^2\cos2\theta$; bij $\cos\theta_{\text{eq}} = g/\omega^2R$ is
dit $m\omega^2R^2\sin^2\theta_{\text{eq}}$, dus
$\omega_{\text{osc}} = \omega\sin\theta_{\text{eq}} = \omega\sqrt{1 -
(g/\omega^2R)^2} \to 0$ als $\omega^2 \to g/R$: de terugdrijvende kracht
vervlakt precies wanneer de putten samenvloeien.
\end{solution}

\begin{solution}{exo:b3:lagrangian-mechanics:10}
(a) Posities $x_2 = \ell(\sin\theta_1 + \sin\theta_2)$ enzovoort; houden
we de kwadratische termen, dan is de gemengde snelheidsterm
$m\ell^2\dot\theta_1\dot\theta_2$, wat de vermelde $L$ geeft. (b)
$2\ddot\theta_1 + \ddot\theta_2 = -2\omega_0^2\theta_1$ en
$\ddot\theta_1 + \ddot\theta_2 = -\omega_0^2\theta_2$, met $\omega_0^2
= g/\ell$. (c) Invullen van $\theta_i = a_i\cos\Omega t$:
$\det\begin{pmatrix} 2\omega_0^2 - 2\Omega^2 & -\Omega^2\\ -\Omega^2 &
\omega_0^2 - \Omega^2\end{pmatrix} = \Omega^4 - 4\omega_0^2\Omega^2 +
2\omega_0^4 = 0$, dus $\Omega_\pm^2 = (2 \mp \sqrt2)\omega_0^2$; de
tweede regel geeft $a_2/a_1 = \Omega^2/(\omega_0^2 - \Omega^2) =
\pm\sqrt2$: de gewichten samen (trage modus), de gewichten tegengesteld
(snelle modus). (d) Bij grote amplitude maken de koppelingen in $\sin$
en $\cos$ de vergelijkingen niet-lineair; oplossingen zijn niet meer
superponeerbaar, en naburige begincondities lopen exponentieel uiteen
(chaos). De energie blijft nog steeds exact behouden --- $L$ bevat de
tijd niet expliciet --- chaos gaat over voorspelbaarheid, niet over
behoud.
\end{solution}

\begin{solution}{exo:b3:lagrangian-mechanics:11}
(a) $v = \sqrt{2gy}$ en $\dd s = \sqrt{1 + y'^2}\,\dd x$ leveren de
functionaal. (b) De berekening van
\cref{prop:b3:lagrangian-mechanics:energy} met $x$ in de rol van de
tijd: $\dd
h/\dd x = -\partial F/\partial x = 0$. (c) $h = -1/\sqrt{2gy(1 +
y'^2)}$, dus $y(1 + y'^2) = 2r$. Voor de cycloïde is $y' =
\sin\phi/(1 - \cos\phi)$ en $1 + y'^2 = 2/(1 - \cos\phi)$, dus
$y(1 + y'^2) = 2r$. (d) $\dd t = \dd s/v = \sqrt{r/g}\,\dd\phi$ (alle
afhankelijkheid van $\phi$ valt weg), dus wordt het laagste punt ($\phi = \pi$) bereikt in $\pi\sqrt{r/g}$ --- vanaf welk beginpunt ook: de
cycloïde is tevens de tautochroon. De rechte glijbaan naar $(\pi r, 2r)$
kost $\sqrt{\pi^2
+ 4}\,\sqrt{r/g} \approx 3.72\sqrt{r/g}$, zo'n $18\%$
meer dan $\pi\sqrt{r/g}$.
\end{solution}

\begin{solution}{exo:b3:lagrangian-mechanics:12}
(a) $\dd t = \dd s/(c/n)$. (b) $h = -n(y)/\sqrt{1 + y'^2}$ blijft
behouden; $1/\sqrt{1 + y'^2} = \cos\theta$ (met $\theta$ de hellingshoek)
$= \sin i$ voor $i$ vanaf de verticaal: $n\sin i = \text{const}$ --- de
wet van Snellius, continu toegepast. (c) Voor een bijna horizontale
straal geeft $n\cos\theta \approx \text{const}$ met kleine $\theta$ dat
$\theta\,\dd\theta = \dd n/n \approx \beta\,\dd y$; omdat $\dd y =
\theta\,\dd x$, is de kromming $\dd\theta/\dd x = \beta$, dus $R =
1/\beta \approx \qty{83}{km}$, met een buiging \emph{omhoog} (naar
grotere $n$). (d) Een straal vanuit het oog die na die kromming met
straal $R$ rakend aan het wegdek loopt, raakt het op $d = \sqrt{2hR} = \sqrt{2 \times 1.2 \times
8.3e4} \approx \qty{450}{m}$: verder dan die
afstand ziet men het wegdek zelf niet --- men ziet omhoog gebroken
hemel, de glinsterende ``waterplas''.
\end{solution}

\begin{solution}{pb:b3:lagrangian-mechanics:1}
\textbf{1.} $L = \tfrac12 m\ell^2\dot\theta^2 + mg\ell\cos\theta$;
$\ddot\theta = -(g/\ell)\sin\theta$.
\textbf{2.} $f_0 = \sqrt{g/\ell}/2\pi = \qty{0.79}{Hz}$.
\textbf{3.} De ophanging staat vast, dus $h = E = \tfrac12
m\ell^2\dot\theta^2 - mg\ell\cos\theta$; van onder naar boven geldt
$\tfrac12
mv^2 = 2mg\ell$: $v = 2\sqrt{g\ell} = \qty{4.0}{m/s}$.
\textbf{4.} $\ddot\epsilon = +(g/\ell)\epsilon$: $\epsilon \propto
\eu^{t/\tau}$ met $\tau = \sqrt{\ell/g} = \qty{0.20}{s}$.
\textbf{5.} Het omgekeerde evenwicht bestaat, maar is exponentieel
onstabiel: elke helling, hoe klein ook, vermenigvuldigt zich elke
\qty{0.2}{s} met $e$.
\textbf{6.} Gewicht in $(\ell\sin\theta,\ a\cos\Omega t - \ell\cos\theta)$; differentiëren en kwadrateren geeft de vermelde
$\vect v^{\,2}$, met de gemengde term
$-2a\Omega\ell\sin(\Omega t)\sin\theta\,\dot\theta$.
\textbf{7.} Is $L' = L + \dd F(q,t)/\dd t$, dan verandert de \omterm{def:b3:lagrangian-mechanics:action}{actie} met
$F(q_2, t_2) - F(q_1, t_1)$, een constante onder variaties met vaste
eindpunten: dezelfde stationaire paden, dezelfde vergelijkingen.
\textbf{8.} De gemengde term is $-ma\Omega\ell\sin(\Omega t)\sin\theta\,
\dot\theta = \dd[ma\Omega\ell\sin(\Omega t)\cos\theta]/\dd t -
ma\Omega^2\ell\cos(\Omega t)\cos\theta$; laten we de totale afgeleide en
de termen in $t$ alleen vallen ($\tfrac12 ma^2\Omega^2\sin^2\Omega t$ en
$-mga\cos\Omega t$), dan blijft de vermelde $L$ over.
\textbf{9.} $m\ell^2\ddot\theta = -mg\ell\sin\theta +
ma\Omega^2\ell\cos(\Omega t)\sin\theta$. Omdat $\ddot y_{\text{s}} =
-a\Omega^2\cos\Omega t$, is dit $\ddot\theta =
-[(g + \ddot y_{\text{s}})/\ell]\sin\theta$: in het assenstelsel van het
draaipunt trilt de schijnbare zwaartekracht.
\textbf{10.} $L$ hangt nu expliciet van $t$ af: $h$ blijft niet behouden
en $E$ evenmin --- de schudder voert via het draaipunt energie aan en af.
\textbf{11.} $a/\ell = 0.05$; $\Omega = \qty{251}{rad/s}$ tegen
$\omega_0 = \qty{4.9}{rad/s}$: verhouding $51$. In één aandrijfperiode
(\qty{25}{ms}) verandert de zwaartekracht $\dot\theta$ nauwelijks: het
gewicht is te traag om te volgen en trilt alleen mee.
\textbf{12.} $\ddot\xi$ is de grootste afgeleide ($\propto\Omega^2$) en
de aandrijving de grootste krachtterm: $\ddot\xi =
(a\Omega^2/\ell)\cos(\Omega t)\sin\Theta$.
\textbf{13.} Tweemaal integreren bij vaste $\Theta$: $\xi =
-(a/\ell)\cos(\Omega t)\sin\Theta$, met een amplitude van hoogstens
$a/\ell =
0.05$: een trilling van twee graden.
\textbf{14.} $\sin\theta \approx \sin\Theta + \xi\cos\Theta$, dus
\[
\ddot\Theta + \ddot\xi = -\frac{g}{\ell}(\sin\Theta + \xi\cos\Theta)
 + \frac{a\Omega^2}{\ell}\cos\Omega t\,(\sin\Theta + \xi\cos\Theta) .
\]
\textbf{15.} Middelen doodt $\ddot\xi$, $\langle\cos\Omega t\rangle$ en
$\langle\xi\rangle$; de overlevende gemengde term is
$(a\Omega^2/\ell)\cos\Theta\,\langle\xi\cos\Omega t\rangle =
-(a^2\Omega^2/2\ell^2)\sin\Theta\cos\Theta$, wat de vermelde
vergelijking voor $\Theta$ geeft.
\textbf{16.} $\ddot\Theta = -(1/m\ell^2)\,U_{\text{eff}}'(\Theta)$ met
$U_{\text{eff}} = mg\ell[-\cos\Theta + (a^2\Omega^2/4g\ell)
\sin^2\Theta]$ --- differentieer ter controle.
\textbf{17.} Dezelfde term in $\sin^2$ als bij de hoepel, maar met het
\emph{tegengestelde} teken: de rotatie groef putten op de flanken en kon
de bodem alleen vlakker maken; het verticale schudden maakt de put
onderaan juist stijver en graaft een nieuwe put \emph{bovenaan}.
\textbf{18.} Trage aandrijving ($a^2\Omega^2 < 2g\ell$): minimum bij
$\Theta =
0$, maximum bij $\pi$ --- niets nieuws. Snelle aandrijving:
minima bij $0$ \emph{en} $\pi$, gescheiden door maxima bij
$\cos\Theta^* =
-2g\ell/a^2\Omega^2$; zowel de hangende als de staande
slinger trilt stabiel.
\textbf{19.} Nabij $\pi$, met $\Theta = \pi + \epsilon$:
$\ddot\epsilon = [g/\ell - a^2\Omega^2/2\ell^2]\,\epsilon$; stabiliteit
vergt een negatief haakje: $a^2\Omega^2 > 2g\ell$.
\textbf{20.} $\Omega_{\text{c}} = \sqrt{2g\ell}/a = \qty{140}{rad/s}$:
$f_{\text{c}} = \qty{22}{Hz}$; pieksnelheid $a\Omega_{\text{c}} =
\qty{2.8}{m/s}$, piekversnelling $a\Omega_{\text{c}}^2 =
\qty{392}{m/s^2} \approx 40g$.
\textbf{21.} $\omega_{\text{traag}} = \sqrt{a^2\Omega^2/2\ell^2 -
g/\ell}$; bij $\Omega = 2\Omega_{\text{c}}$ is $a^2\Omega^2 = 8g\ell$,
dus $\omega_{\text{traag}} = \sqrt{3g/\ell} = \qty{8.6}{rad/s}$:
\qty{1.4}{Hz}, dertig keer trager dan de aandrijving van \qty{45}{Hz}.
\textbf{22.} $\xi_{\max} = (a/\ell)\sin\Theta$: bij $10^\circ$ uit de
verticaal is $\xi_{\max} = 0.05\sin170^\circ = \qty{8.7}{mrad}
\approx 0.5^\circ$ --- op een gewone foto staat de bezem doodstil.
\textbf{23.} De staande put reikt tot de flankerende maxima:
$|\Theta - \pi| < \arccos(2g\ell/a^2\Omega^2)$; bij $\Omega =
2\Omega_{\text{c}}$ is $\arccos\tfrac14 = 75^\circ$ --- een opmerkelijk
vergevingsgezinde put.
\textbf{24.} De jongleur gebruikt terugkoppeling: de ogen meten de
helling, de hand versnelt zijwaarts om haar op te heffen. De
stabilisatie van Kapitza werkt in open lus --- de aandrijving weet nooit
waar de slinger staat.
\textbf{25.} Geschud met \qty{45}{Hz} en een amplitude van \qty{2}{cm}
staat de slinger van \qty{40}{cm} omgekeerd in een put van $75^\circ$ en
wiegt hij met \qty{1.4}{Hz}. Dezelfde uitgemiddelde effectieve
potentiaal, gemaakt met een oscillerend elektrisch quadrupoolveld in
plaats van een geschud draaipunt, sluit losse ionen op in de val van
Paul --- het werkpaard van atoomklokken en van kwantumcomputers met
gevangen ionen.
\end{solution}
